%% =============================================================================
%% section6.tex — Section 6: Discussion
%% =============================================================================

\section{Discussion}\label{sec:discussion}

%% ---------------------------------------------------------------------------
%% Summary of contributions (methodology-first framing)
%% ---------------------------------------------------------------------------

\paragraph{Summary.}
This paper develops a hierarchical hurdle beta-binomial (HBB) framework
for survey-weighted bounded counts with structural zeros, extending the
pseudo-posterior methodology of~\citet{WilliamsSavitsky2021} to two-part
hierarchical models.  The two primary methodological
contributions---cross-margin covariance identification
(\cref{prop:sigma12-nonid}) and $\DER$-guided sandwich correction
(\cref{thm:cholesky})---are discussed below alongside generalizability,
practical recommendations, and limitations.

%% ---------------------------------------------------------------------------
%% Methodological contributions (compressed: 4 domains → 2)
%% ---------------------------------------------------------------------------

\paragraph{Generalizability.}
The HBB framework developed in~\cref{sec:model} applies whenever the
outcome is a bounded discrete proportion---a count
$Y_i \in \{0, 1, \ldots, n_i\}$ with a known denominator---embedded
in a hierarchical structure with structural zeros.  Five domains
illustrate the breadth of applicability beyond childcare enrollment.
(i)~In \emph{dental epidemiology}, $Y_i$ is the number of carious
surfaces out of $n_i$ surfaces examined, with structural zeros for
caries-free individuals whose zero counts reflect disease absence
rather than low risk~\citep{BandyopadhyayEtAl2011}.
(ii)~In \emph{species occupancy modeling}, $Y_i$ is the number of
detection events out of $n_i$ replicate surveys at a site, with
structural zeros at sites where the species is truly
absent~\citep{MacKenzieEtAl2002}.
(iii)~In \emph{hospital bed management}, $Y_i$ is the number of beds
occupied out of $n_i$ staffed beds on a ward, with structural zeros
for wards temporarily closed to admissions due to staffing shortages
or infection control measures.
(iv)~In \emph{insurance enrollment}, $Y_i$ is the number of enrollees
out of $n_i$ eligible individuals within a plan--region cell, with
structural zeros for plan--region combinations in which the plan is
not offered.
(v)~In \emph{survey item endorsement}, $Y_i$ is the number of items
endorsed out of $n_i$ items administered on a subscale, with
structural zeros for respondents to whom the construct does not apply
(e.g., substance-use items for lifetime abstainers).
Each domain exhibits the three features that jointly motivate the HBB
specification: bounded support with a known denominator,
extra-binomial variation within the active population, and a mass of
structural zeros generated by a qualitatively distinct mechanism from
the count process.  In every case, the hurdle separates the structural
participation decision from the intensity conditional on
participation, and the beta-binomial kernel accommodates
overdispersion relative to the binomial baseline.

\paragraph{Inference contributions.}
Beyond the modeling framework, this paper makes two specific
contributions to inference.  First,
\cref{prop:sigma12-nonid} establishes that the observation-level
conditional likelihood is block-diagonal across margins, so the
cross-margin covariance $\bSigma_{12}$ (\cref{sec:svc}) is a
hierarchical estimand whose effective sample size is~$S$, not~$N$.
While a frequentist could in principle target $\bSigma_{12}$ via REML
on the marginal likelihood (\cref{rem:cond-vs-marg}), unrestricted estimation of a
$q \times q$ cross-covariance from $S = 51$ groups is fragile when
$2q = 10$; \cref{prop:sigma12-id} shows that the $\LKJ$ prior provides
finite-sample regularization, yielding posterior concentration at rate
$O(S^{-1/2})$.
Second, the sandwich correction pipeline with $\DER$-guided block-wise
application (\cref{thm:cholesky}) extends the pseudo-posterior
framework of~\citet{WilliamsSavitsky2021} to two-part hierarchical
models, nesting the scalar design-effect multiplier
of~\citet{GhosalEtAl2020} as a special case when the $\DER$ is
constant across parameters (\cref{eq:der}).

%% ---------------------------------------------------------------------------
%% Practical recommendations (6 → 3)
%% ---------------------------------------------------------------------------

\paragraph{Practical recommendations.}
The combined evidence from the application (\cref{sec:application})
and simulation (\cref{sec:simulation}) suggests three guidelines for
applied researchers fitting Bayesian hierarchical models to complex
survey data.
First, \emph{test for informativeness before deciding on weighting}.
Following~\citet{Pfeffermann1993}, regress the sampling weight on the
outcome conditional on the full model covariate vector and test whether
the weight coefficient is significant; conditioning on the full
covariate set is essential, because a univariate test conflates
informativeness with confounding~\citep{PfeffermannEtAl1998}. For
two-part models, test each margin separately, as the design can be
non-informative for one margin and informative for the other
(\cref{sec:data}).

Second, \emph{report both naive and sandwich-corrected confidence
intervals alongside the $\DER$}.
The empirical $\DER$ range in this study is 1.14--4.18
(\cref{tab:fixed-effects}), and the simulation width ratios of
1.22--2.06 (\cref{tab:sim-results}) provide independent validation.
The linear interpolation heuristic overestimates the $\DER$ at
intermediate shrinkage levels by up to 50\%; the exact conjugate
formula in Supplementary Material~C (\cref{app:survey}) provides a
more accurate benchmark.

Third, \emph{for hierarchical variance components, report model-based
posterior intervals and conduct prior sensitivity analysis}.
The sandwich correction is structurally inapplicable to
parameters---such as the state-effect standard deviations $\tau$ and
the cross-margin correlation $\rho$---that do not produce
individual-level score contributions (\cref{sec:simresults}), and the
simulation confirms that the width ratio is identically 1.00 for these
parameters. We conduct a sensitivity analysis in
\cref{smd:lkj-sensitivity}, refitting under $\LKJ(\eta)$ with
$\eta \in \{1,\ldots,8\}$: the cross-margin correlation remains near
zero ($\hat{\rho} \in [0.10,\,0.19]$) and all fixed-effect estimates
are stable, confirming that the posterior is data-dominated.
For the 3--5 largest clusters, a sensitivity analysis comparing
block-wise and joint Cholesky correction is also recommended
(\cref{rem:blockwise}; Supplementary Material~C).

%% ---------------------------------------------------------------------------
%% Limitations (6 → 4, with future directions merged)
%% ---------------------------------------------------------------------------

\paragraph{Limitations and future directions.}
Several limitations merit discussion.
The weight ratio $w_{\max}/w_{\min} = 462$ exceeds
$N^{1/2} \approx 82.4$ (\cref{sec:survey}), so the sufficient
condition~C3 for the Bernstein--von Mises theorem under informative
sampling is not comfortably satisfied; this motivates the sandwich
correction as a finite-sample robustness measure rather than an
asymptotic guarantee, and underscores the importance of the
simulation study in~\cref{sec:simulation} as an independent
check on coverage properties.
The beta-binomial component conditions on $n_i$ as a known upper
bound, yet total enrollment may itself respond to community
characteristics---a form of outcome endogeneity.  Joint modeling of
$(Y_i, n_i)$ would address this concern but requires a simultaneous
specification for the total-enrollment and IT-share processes, an
extension we leave to future work.
The model assumes a single concentration parameter $\kappa$ shared
across all states; state-varying $\kappa_s$ could capture
heterogeneous overdispersion, but identification of each $\kappa_s$
requires sufficient within-state replication of the denominator
values~$n_i$ (\cref{smb:prop-mu-kappa-id}).
Finally, the sandwich correction delivers design-consistent inference
for fixed effects but does not extend to hierarchical variance
components such as $\bSigma_{12}$; extending design-consistent
inference to these parameters remains an open problem, with
replicate-weight approaches offering a potential but computationally
costly avenue.

\paragraph{Alternative survey approaches.}
The pseudo-posterior framework adopted here is not the only route to
survey-weighted Bayesian inference.  Multilevel pseudo-maximum
likelihood~\citep[MPML;][]{RabeHeskethSkrondal2006} provides a
frequentist counterpart that integrates sampling weights directly into
the marginal likelihood, avoiding the pseudo-likelihood construction.
Fully model-based approaches that include the sampling mechanism as part
of the probability model~\citep{SiEtAl2020} offer theoretical
advantages when the design is informative but require additional
specification of the selection process.  Finally, balanced repeated
replicate (BRR) or jackknife replicate-weight
methods~\citep{RaoWu1988} provide design-based variance estimates
without distributional assumptions, though the computational cost of
refitting Bayesian hierarchical models under each set of replicate
weights is currently prohibitive.  Comparing these alternatives with
the pseudo-posterior sandwich approach is an important direction for
future work.

\paragraph{Conclusion.}
This paper introduces a Bayesian hierarchical hurdle beta-binomial
framework for survey-weighted bounded counts with structural zeros.
Applied to the 2019 NSECE, the model reveals a poverty reversal---community
poverty reduces access to infant and toddler care but increases
enrollment intensity among participating centers---that is
near-universal across states under hierarchical shrinkage (48 of 51)
and whose extensive margin accounts for approximately two-thirds of
the total effect.  The sandwich correction and DER diagnostic provide
a practical toolkit for population-level inference in Bayesian
hierarchical models fitted to complex survey data.  The companion R
package \textsf{hurdlebb}~\citep{hurdlebb2026} and replication
repository make the methodology immediately accessible to applied
researchers.

%% End-matter statements are in main.tex (after %% =============================================================================
%% section6.tex — Section 6: Discussion
%% =============================================================================

\section{Discussion}\label{sec:discussion}

%% ---------------------------------------------------------------------------
%% Summary of contributions (methodology-first framing)
%% ---------------------------------------------------------------------------

\paragraph{Summary.}
This paper develops a hierarchical hurdle beta-binomial (HBB) framework
for survey-weighted bounded counts with structural zeros, extending the
pseudo-posterior methodology of~\citet{WilliamsSavitsky2021} to two-part
hierarchical models.  The two primary methodological
contributions---cross-margin covariance identification
(\cref{prop:sigma12-nonid}) and $\DER$-guided sandwich correction
(\cref{thm:cholesky})---are discussed below alongside generalizability,
practical recommendations, and limitations.

%% ---------------------------------------------------------------------------
%% Methodological contributions (compressed: 4 domains → 2)
%% ---------------------------------------------------------------------------

\paragraph{Generalizability.}
The HBB framework developed in~\cref{sec:model} applies whenever the
outcome is a bounded discrete proportion---a count
$Y_i \in \{0, 1, \ldots, n_i\}$ with a known denominator---embedded
in a hierarchical structure with structural zeros.  Five domains
illustrate the breadth of applicability beyond childcare enrollment.
(i)~In \emph{dental epidemiology}, $Y_i$ is the number of carious
surfaces out of $n_i$ surfaces examined, with structural zeros for
caries-free individuals whose zero counts reflect disease absence
rather than low risk~\citep{BandyopadhyayEtAl2011}.
(ii)~In \emph{species occupancy modeling}, $Y_i$ is the number of
detection events out of $n_i$ replicate surveys at a site, with
structural zeros at sites where the species is truly
absent~\citep{MacKenzieEtAl2002}.
(iii)~In \emph{hospital bed management}, $Y_i$ is the number of beds
occupied out of $n_i$ staffed beds on a ward, with structural zeros
for wards temporarily closed to admissions due to staffing shortages
or infection control measures.
(iv)~In \emph{insurance enrollment}, $Y_i$ is the number of enrollees
out of $n_i$ eligible individuals within a plan--region cell, with
structural zeros for plan--region combinations in which the plan is
not offered.
(v)~In \emph{survey item endorsement}, $Y_i$ is the number of items
endorsed out of $n_i$ items administered on a subscale, with
structural zeros for respondents to whom the construct does not apply
(e.g., substance-use items for lifetime abstainers).
Each domain exhibits the three features that jointly motivate the HBB
specification: bounded support with a known denominator,
extra-binomial variation within the active population, and a mass of
structural zeros generated by a qualitatively distinct mechanism from
the count process.  In every case, the hurdle separates the structural
participation decision from the intensity conditional on
participation, and the beta-binomial kernel accommodates
overdispersion relative to the binomial baseline.

\paragraph{Inference contributions.}
Beyond the modeling framework, this paper makes two specific
contributions to inference.  First,
\cref{prop:sigma12-nonid} establishes that the observation-level
conditional likelihood is block-diagonal across margins, so the
cross-margin covariance $\bSigma_{12}$ (\cref{sec:svc}) is a
hierarchical estimand whose effective sample size is~$S$, not~$N$.
While a frequentist could in principle target $\bSigma_{12}$ via REML
on the marginal likelihood (\cref{rem:cond-vs-marg}), unrestricted estimation of a
$q \times q$ cross-covariance from $S = 51$ groups is fragile when
$2q = 10$; \cref{prop:sigma12-id} shows that the $\LKJ$ prior provides
finite-sample regularization, yielding posterior concentration at rate
$O(S^{-1/2})$.
Second, the sandwich correction pipeline with $\DER$-guided block-wise
application (\cref{thm:cholesky}) extends the pseudo-posterior
framework of~\citet{WilliamsSavitsky2021} to two-part hierarchical
models, nesting the scalar design-effect multiplier
of~\citet{GhosalEtAl2020} as a special case when the $\DER$ is
constant across parameters (\cref{eq:der}).

%% ---------------------------------------------------------------------------
%% Practical recommendations (6 → 3)
%% ---------------------------------------------------------------------------

\paragraph{Practical recommendations.}
The combined evidence from the application (\cref{sec:application})
and simulation (\cref{sec:simulation}) suggests three guidelines for
applied researchers fitting Bayesian hierarchical models to complex
survey data.
First, \emph{test for informativeness before deciding on weighting}.
Following~\citet{Pfeffermann1993}, regress the sampling weight on the
outcome conditional on the full model covariate vector and test whether
the weight coefficient is significant; conditioning on the full
covariate set is essential, because a univariate test conflates
informativeness with confounding~\citep{PfeffermannEtAl1998}. For
two-part models, test each margin separately, as the design can be
non-informative for one margin and informative for the other
(\cref{sec:data}).

Second, \emph{report both naive and sandwich-corrected confidence
intervals alongside the $\DER$}.
The empirical $\DER$ range in this study is 1.14--4.18
(\cref{tab:fixed-effects}), and the simulation width ratios of
1.22--2.06 (\cref{tab:sim-results}) provide independent validation.
The linear interpolation heuristic overestimates the $\DER$ at
intermediate shrinkage levels by up to 50\%; the exact conjugate
formula in Supplementary Material~C (\cref{app:survey}) provides a
more accurate benchmark.

Third, \emph{for hierarchical variance components, report model-based
posterior intervals and conduct prior sensitivity analysis}.
The sandwich correction is structurally inapplicable to
parameters---such as the state-effect standard deviations $\tau$ and
the cross-margin correlation $\rho$---that do not produce
individual-level score contributions (\cref{sec:simresults}), and the
simulation confirms that the width ratio is identically 1.00 for these
parameters. We conduct a sensitivity analysis in
\cref{smd:lkj-sensitivity}, refitting under $\LKJ(\eta)$ with
$\eta \in \{1,\ldots,8\}$: the cross-margin correlation remains near
zero ($\hat{\rho} \in [0.10,\,0.19]$) and all fixed-effect estimates
are stable, confirming that the posterior is data-dominated.
For the 3--5 largest clusters, a sensitivity analysis comparing
block-wise and joint Cholesky correction is also recommended
(\cref{rem:blockwise}; Supplementary Material~C).

%% ---------------------------------------------------------------------------
%% Limitations (6 → 4, with future directions merged)
%% ---------------------------------------------------------------------------

\paragraph{Limitations and future directions.}
Several limitations merit discussion.
The weight ratio $w_{\max}/w_{\min} = 462$ exceeds
$N^{1/2} \approx 82.4$ (\cref{sec:survey}), so the sufficient
condition~C3 for the Bernstein--von Mises theorem under informative
sampling is not comfortably satisfied; this motivates the sandwich
correction as a finite-sample robustness measure rather than an
asymptotic guarantee, and underscores the importance of the
simulation study in~\cref{sec:simulation} as an independent
check on coverage properties.
The beta-binomial component conditions on $n_i$ as a known upper
bound, yet total enrollment may itself respond to community
characteristics---a form of outcome endogeneity.  Joint modeling of
$(Y_i, n_i)$ would address this concern but requires a simultaneous
specification for the total-enrollment and IT-share processes, an
extension we leave to future work.
The model assumes a single concentration parameter $\kappa$ shared
across all states; state-varying $\kappa_s$ could capture
heterogeneous overdispersion, but identification of each $\kappa_s$
requires sufficient within-state replication of the denominator
values~$n_i$ (\cref{smb:prop-mu-kappa-id}).
Finally, the sandwich correction delivers design-consistent inference
for fixed effects but does not extend to hierarchical variance
components such as $\bSigma_{12}$; extending design-consistent
inference to these parameters remains an open problem, with
replicate-weight approaches offering a potential but computationally
costly avenue.

\paragraph{Alternative survey approaches.}
The pseudo-posterior framework adopted here is not the only route to
survey-weighted Bayesian inference.  Multilevel pseudo-maximum
likelihood~\citep[MPML;][]{RabeHeskethSkrondal2006} provides a
frequentist counterpart that integrates sampling weights directly into
the marginal likelihood, avoiding the pseudo-likelihood construction.
Fully model-based approaches that include the sampling mechanism as part
of the probability model~\citep{SiEtAl2020} offer theoretical
advantages when the design is informative but require additional
specification of the selection process.  Finally, balanced repeated
replicate (BRR) or jackknife replicate-weight
methods~\citep{RaoWu1988} provide design-based variance estimates
without distributional assumptions, though the computational cost of
refitting Bayesian hierarchical models under each set of replicate
weights is currently prohibitive.  Comparing these alternatives with
the pseudo-posterior sandwich approach is an important direction for
future work.

\paragraph{Conclusion.}
This paper introduces a Bayesian hierarchical hurdle beta-binomial
framework for survey-weighted bounded counts with structural zeros.
Applied to the 2019 NSECE, the model reveals a poverty reversal---community
poverty reduces access to infant and toddler care but increases
enrollment intensity among participating centers---that is
near-universal across states under hierarchical shrinkage (48 of 51)
and whose extensive margin accounts for approximately two-thirds of
the total effect.  The sandwich correction and DER diagnostic provide
a practical toolkit for population-level inference in Bayesian
hierarchical models fitted to complex survey data.  The companion R
package \textsf{hurdlebb}~\citep{hurdlebb2026} and replication
repository make the methodology immediately accessible to applied
researchers.

%% End-matter statements are in main.tex (after %% =============================================================================
%% section6.tex — Section 6: Discussion
%% =============================================================================

\section{Discussion}\label{sec:discussion}

%% ---------------------------------------------------------------------------
%% Summary of contributions (methodology-first framing)
%% ---------------------------------------------------------------------------

\paragraph{Summary.}
This paper develops a hierarchical hurdle beta-binomial (HBB) framework
for survey-weighted bounded counts with structural zeros, extending the
pseudo-posterior methodology of~\citet{WilliamsSavitsky2021} to two-part
hierarchical models.  The two primary methodological
contributions---cross-margin covariance identification
(\cref{prop:sigma12-nonid}) and $\DER$-guided sandwich correction
(\cref{thm:cholesky})---are discussed below alongside generalizability,
practical recommendations, and limitations.

%% ---------------------------------------------------------------------------
%% Methodological contributions (compressed: 4 domains → 2)
%% ---------------------------------------------------------------------------

\paragraph{Generalizability.}
The HBB framework developed in~\cref{sec:model} applies whenever the
outcome is a bounded discrete proportion---a count
$Y_i \in \{0, 1, \ldots, n_i\}$ with a known denominator---embedded
in a hierarchical structure with structural zeros.  Five domains
illustrate the breadth of applicability beyond childcare enrollment.
(i)~In \emph{dental epidemiology}, $Y_i$ is the number of carious
surfaces out of $n_i$ surfaces examined, with structural zeros for
caries-free individuals whose zero counts reflect disease absence
rather than low risk~\citep{BandyopadhyayEtAl2011}.
(ii)~In \emph{species occupancy modeling}, $Y_i$ is the number of
detection events out of $n_i$ replicate surveys at a site, with
structural zeros at sites where the species is truly
absent~\citep{MacKenzieEtAl2002}.
(iii)~In \emph{hospital bed management}, $Y_i$ is the number of beds
occupied out of $n_i$ staffed beds on a ward, with structural zeros
for wards temporarily closed to admissions due to staffing shortages
or infection control measures.
(iv)~In \emph{insurance enrollment}, $Y_i$ is the number of enrollees
out of $n_i$ eligible individuals within a plan--region cell, with
structural zeros for plan--region combinations in which the plan is
not offered.
(v)~In \emph{survey item endorsement}, $Y_i$ is the number of items
endorsed out of $n_i$ items administered on a subscale, with
structural zeros for respondents to whom the construct does not apply
(e.g., substance-use items for lifetime abstainers).
Each domain exhibits the three features that jointly motivate the HBB
specification: bounded support with a known denominator,
extra-binomial variation within the active population, and a mass of
structural zeros generated by a qualitatively distinct mechanism from
the count process.  In every case, the hurdle separates the structural
participation decision from the intensity conditional on
participation, and the beta-binomial kernel accommodates
overdispersion relative to the binomial baseline.

\paragraph{Inference contributions.}
Beyond the modeling framework, this paper makes two specific
contributions to inference.  First,
\cref{prop:sigma12-nonid} establishes that the observation-level
conditional likelihood is block-diagonal across margins, so the
cross-margin covariance $\bSigma_{12}$ (\cref{sec:svc}) is a
hierarchical estimand whose effective sample size is~$S$, not~$N$.
While a frequentist could in principle target $\bSigma_{12}$ via REML
on the marginal likelihood (\cref{rem:cond-vs-marg}), unrestricted estimation of a
$q \times q$ cross-covariance from $S = 51$ groups is fragile when
$2q = 10$; \cref{prop:sigma12-id} shows that the $\LKJ$ prior provides
finite-sample regularization, yielding posterior concentration at rate
$O(S^{-1/2})$.
Second, the sandwich correction pipeline with $\DER$-guided block-wise
application (\cref{thm:cholesky}) extends the pseudo-posterior
framework of~\citet{WilliamsSavitsky2021} to two-part hierarchical
models, nesting the scalar design-effect multiplier
of~\citet{GhosalEtAl2020} as a special case when the $\DER$ is
constant across parameters (\cref{eq:der}).

%% ---------------------------------------------------------------------------
%% Practical recommendations (6 → 3)
%% ---------------------------------------------------------------------------

\paragraph{Practical recommendations.}
The combined evidence from the application (\cref{sec:application})
and simulation (\cref{sec:simulation}) suggests three guidelines for
applied researchers fitting Bayesian hierarchical models to complex
survey data.
First, \emph{test for informativeness before deciding on weighting}.
Following~\citet{Pfeffermann1993}, regress the sampling weight on the
outcome conditional on the full model covariate vector and test whether
the weight coefficient is significant; conditioning on the full
covariate set is essential, because a univariate test conflates
informativeness with confounding~\citep{PfeffermannEtAl1998}. For
two-part models, test each margin separately, as the design can be
non-informative for one margin and informative for the other
(\cref{sec:data}).

Second, \emph{report both naive and sandwich-corrected confidence
intervals alongside the $\DER$}.
The empirical $\DER$ range in this study is 1.14--4.18
(\cref{tab:fixed-effects}), and the simulation width ratios of
1.22--2.06 (\cref{tab:sim-results}) provide independent validation.
The linear interpolation heuristic overestimates the $\DER$ at
intermediate shrinkage levels by up to 50\%; the exact conjugate
formula in Supplementary Material~C (\cref{app:survey}) provides a
more accurate benchmark.

Third, \emph{for hierarchical variance components, report model-based
posterior intervals and conduct prior sensitivity analysis}.
The sandwich correction is structurally inapplicable to
parameters---such as the state-effect standard deviations $\tau$ and
the cross-margin correlation $\rho$---that do not produce
individual-level score contributions (\cref{sec:simresults}), and the
simulation confirms that the width ratio is identically 1.00 for these
parameters. We conduct a sensitivity analysis in
\cref{smd:lkj-sensitivity}, refitting under $\LKJ(\eta)$ with
$\eta \in \{1,\ldots,8\}$: the cross-margin correlation remains near
zero ($\hat{\rho} \in [0.10,\,0.19]$) and all fixed-effect estimates
are stable, confirming that the posterior is data-dominated.
For the 3--5 largest clusters, a sensitivity analysis comparing
block-wise and joint Cholesky correction is also recommended
(\cref{rem:blockwise}; Supplementary Material~C).

%% ---------------------------------------------------------------------------
%% Limitations (6 → 4, with future directions merged)
%% ---------------------------------------------------------------------------

\paragraph{Limitations and future directions.}
Several limitations merit discussion.
The weight ratio $w_{\max}/w_{\min} = 462$ exceeds
$N^{1/2} \approx 82.4$ (\cref{sec:survey}), so the sufficient
condition~C3 for the Bernstein--von Mises theorem under informative
sampling is not comfortably satisfied; this motivates the sandwich
correction as a finite-sample robustness measure rather than an
asymptotic guarantee, and underscores the importance of the
simulation study in~\cref{sec:simulation} as an independent
check on coverage properties.
The beta-binomial component conditions on $n_i$ as a known upper
bound, yet total enrollment may itself respond to community
characteristics---a form of outcome endogeneity.  Joint modeling of
$(Y_i, n_i)$ would address this concern but requires a simultaneous
specification for the total-enrollment and IT-share processes, an
extension we leave to future work.
The model assumes a single concentration parameter $\kappa$ shared
across all states; state-varying $\kappa_s$ could capture
heterogeneous overdispersion, but identification of each $\kappa_s$
requires sufficient within-state replication of the denominator
values~$n_i$ (\cref{smb:prop-mu-kappa-id}).
Finally, the sandwich correction delivers design-consistent inference
for fixed effects but does not extend to hierarchical variance
components such as $\bSigma_{12}$; extending design-consistent
inference to these parameters remains an open problem, with
replicate-weight approaches offering a potential but computationally
costly avenue.

\paragraph{Alternative survey approaches.}
The pseudo-posterior framework adopted here is not the only route to
survey-weighted Bayesian inference.  Multilevel pseudo-maximum
likelihood~\citep[MPML;][]{RabeHeskethSkrondal2006} provides a
frequentist counterpart that integrates sampling weights directly into
the marginal likelihood, avoiding the pseudo-likelihood construction.
Fully model-based approaches that include the sampling mechanism as part
of the probability model~\citep{SiEtAl2020} offer theoretical
advantages when the design is informative but require additional
specification of the selection process.  Finally, balanced repeated
replicate (BRR) or jackknife replicate-weight
methods~\citep{RaoWu1988} provide design-based variance estimates
without distributional assumptions, though the computational cost of
refitting Bayesian hierarchical models under each set of replicate
weights is currently prohibitive.  Comparing these alternatives with
the pseudo-posterior sandwich approach is an important direction for
future work.

\paragraph{Conclusion.}
This paper introduces a Bayesian hierarchical hurdle beta-binomial
framework for survey-weighted bounded counts with structural zeros.
Applied to the 2019 NSECE, the model reveals a poverty reversal---community
poverty reduces access to infant and toddler care but increases
enrollment intensity among participating centers---that is
near-universal across states under hierarchical shrinkage (48 of 51)
and whose extensive margin accounts for approximately two-thirds of
the total effect.  The sandwich correction and DER diagnostic provide
a practical toolkit for population-level inference in Bayesian
hierarchical models fitted to complex survey data.  The companion R
package \textsf{hurdlebb}~\citep{hurdlebb2026} and replication
repository make the methodology immediately accessible to applied
researchers.

%% End-matter statements are in main.tex (after %% =============================================================================
%% section6.tex — Section 6: Discussion
%% =============================================================================

\section{Discussion}\label{sec:discussion}

%% ---------------------------------------------------------------------------
%% Summary of contributions (methodology-first framing)
%% ---------------------------------------------------------------------------

\paragraph{Summary.}
This paper develops a hierarchical hurdle beta-binomial (HBB) framework
for survey-weighted bounded counts with structural zeros, extending the
pseudo-posterior methodology of~\citet{WilliamsSavitsky2021} to two-part
hierarchical models.  The two primary methodological
contributions---cross-margin covariance identification
(\cref{prop:sigma12-nonid}) and $\DER$-guided sandwich correction
(\cref{thm:cholesky})---are discussed below alongside generalizability,
practical recommendations, and limitations.

%% ---------------------------------------------------------------------------
%% Methodological contributions (compressed: 4 domains → 2)
%% ---------------------------------------------------------------------------

\paragraph{Generalizability.}
The HBB framework developed in~\cref{sec:model} applies whenever the
outcome is a bounded discrete proportion---a count
$Y_i \in \{0, 1, \ldots, n_i\}$ with a known denominator---embedded
in a hierarchical structure with structural zeros.  Five domains
illustrate the breadth of applicability beyond childcare enrollment.
(i)~In \emph{dental epidemiology}, $Y_i$ is the number of carious
surfaces out of $n_i$ surfaces examined, with structural zeros for
caries-free individuals whose zero counts reflect disease absence
rather than low risk~\citep{BandyopadhyayEtAl2011}.
(ii)~In \emph{species occupancy modeling}, $Y_i$ is the number of
detection events out of $n_i$ replicate surveys at a site, with
structural zeros at sites where the species is truly
absent~\citep{MacKenzieEtAl2002}.
(iii)~In \emph{hospital bed management}, $Y_i$ is the number of beds
occupied out of $n_i$ staffed beds on a ward, with structural zeros
for wards temporarily closed to admissions due to staffing shortages
or infection control measures.
(iv)~In \emph{insurance enrollment}, $Y_i$ is the number of enrollees
out of $n_i$ eligible individuals within a plan--region cell, with
structural zeros for plan--region combinations in which the plan is
not offered.
(v)~In \emph{survey item endorsement}, $Y_i$ is the number of items
endorsed out of $n_i$ items administered on a subscale, with
structural zeros for respondents to whom the construct does not apply
(e.g., substance-use items for lifetime abstainers).
Each domain exhibits the three features that jointly motivate the HBB
specification: bounded support with a known denominator,
extra-binomial variation within the active population, and a mass of
structural zeros generated by a qualitatively distinct mechanism from
the count process.  In every case, the hurdle separates the structural
participation decision from the intensity conditional on
participation, and the beta-binomial kernel accommodates
overdispersion relative to the binomial baseline.

\paragraph{Inference contributions.}
Beyond the modeling framework, this paper makes two specific
contributions to inference.  First,
\cref{prop:sigma12-nonid} establishes that the observation-level
conditional likelihood is block-diagonal across margins, so the
cross-margin covariance $\bSigma_{12}$ (\cref{sec:svc}) is a
hierarchical estimand whose effective sample size is~$S$, not~$N$.
While a frequentist could in principle target $\bSigma_{12}$ via REML
on the marginal likelihood (\cref{rem:cond-vs-marg}), unrestricted estimation of a
$q \times q$ cross-covariance from $S = 51$ groups is fragile when
$2q = 10$; \cref{prop:sigma12-id} shows that the $\LKJ$ prior provides
finite-sample regularization, yielding posterior concentration at rate
$O(S^{-1/2})$.
Second, the sandwich correction pipeline with $\DER$-guided block-wise
application (\cref{thm:cholesky}) extends the pseudo-posterior
framework of~\citet{WilliamsSavitsky2021} to two-part hierarchical
models, nesting the scalar design-effect multiplier
of~\citet{GhosalEtAl2020} as a special case when the $\DER$ is
constant across parameters (\cref{eq:der}).

%% ---------------------------------------------------------------------------
%% Practical recommendations (6 → 3)
%% ---------------------------------------------------------------------------

\paragraph{Practical recommendations.}
The combined evidence from the application (\cref{sec:application})
and simulation (\cref{sec:simulation}) suggests three guidelines for
applied researchers fitting Bayesian hierarchical models to complex
survey data.
First, \emph{test for informativeness before deciding on weighting}.
Following~\citet{Pfeffermann1993}, regress the sampling weight on the
outcome conditional on the full model covariate vector and test whether
the weight coefficient is significant; conditioning on the full
covariate set is essential, because a univariate test conflates
informativeness with confounding~\citep{PfeffermannEtAl1998}. For
two-part models, test each margin separately, as the design can be
non-informative for one margin and informative for the other
(\cref{sec:data}).

Second, \emph{report both naive and sandwich-corrected confidence
intervals alongside the $\DER$}.
The empirical $\DER$ range in this study is 1.14--4.18
(\cref{tab:fixed-effects}), and the simulation width ratios of
1.22--2.06 (\cref{tab:sim-results}) provide independent validation.
The linear interpolation heuristic overestimates the $\DER$ at
intermediate shrinkage levels by up to 50\%; the exact conjugate
formula in Supplementary Material~C (\cref{app:survey}) provides a
more accurate benchmark.

Third, \emph{for hierarchical variance components, report model-based
posterior intervals and conduct prior sensitivity analysis}.
The sandwich correction is structurally inapplicable to
parameters---such as the state-effect standard deviations $\tau$ and
the cross-margin correlation $\rho$---that do not produce
individual-level score contributions (\cref{sec:simresults}), and the
simulation confirms that the width ratio is identically 1.00 for these
parameters. We conduct a sensitivity analysis in
\cref{smd:lkj-sensitivity}, refitting under $\LKJ(\eta)$ with
$\eta \in \{1,\ldots,8\}$: the cross-margin correlation remains near
zero ($\hat{\rho} \in [0.10,\,0.19]$) and all fixed-effect estimates
are stable, confirming that the posterior is data-dominated.
For the 3--5 largest clusters, a sensitivity analysis comparing
block-wise and joint Cholesky correction is also recommended
(\cref{rem:blockwise}; Supplementary Material~C).

%% ---------------------------------------------------------------------------
%% Limitations (6 → 4, with future directions merged)
%% ---------------------------------------------------------------------------

\paragraph{Limitations and future directions.}
Several limitations merit discussion.
The weight ratio $w_{\max}/w_{\min} = 462$ exceeds
$N^{1/2} \approx 82.4$ (\cref{sec:survey}), so the sufficient
condition~C3 for the Bernstein--von Mises theorem under informative
sampling is not comfortably satisfied; this motivates the sandwich
correction as a finite-sample robustness measure rather than an
asymptotic guarantee, and underscores the importance of the
simulation study in~\cref{sec:simulation} as an independent
check on coverage properties.
The beta-binomial component conditions on $n_i$ as a known upper
bound, yet total enrollment may itself respond to community
characteristics---a form of outcome endogeneity.  Joint modeling of
$(Y_i, n_i)$ would address this concern but requires a simultaneous
specification for the total-enrollment and IT-share processes, an
extension we leave to future work.
The model assumes a single concentration parameter $\kappa$ shared
across all states; state-varying $\kappa_s$ could capture
heterogeneous overdispersion, but identification of each $\kappa_s$
requires sufficient within-state replication of the denominator
values~$n_i$ (\cref{smb:prop-mu-kappa-id}).
Finally, the sandwich correction delivers design-consistent inference
for fixed effects but does not extend to hierarchical variance
components such as $\bSigma_{12}$; extending design-consistent
inference to these parameters remains an open problem, with
replicate-weight approaches offering a potential but computationally
costly avenue.

\paragraph{Alternative survey approaches.}
The pseudo-posterior framework adopted here is not the only route to
survey-weighted Bayesian inference.  Multilevel pseudo-maximum
likelihood~\citep[MPML;][]{RabeHeskethSkrondal2006} provides a
frequentist counterpart that integrates sampling weights directly into
the marginal likelihood, avoiding the pseudo-likelihood construction.
Fully model-based approaches that include the sampling mechanism as part
of the probability model~\citep{SiEtAl2020} offer theoretical
advantages when the design is informative but require additional
specification of the selection process.  Finally, balanced repeated
replicate (BRR) or jackknife replicate-weight
methods~\citep{RaoWu1988} provide design-based variance estimates
without distributional assumptions, though the computational cost of
refitting Bayesian hierarchical models under each set of replicate
weights is currently prohibitive.  Comparing these alternatives with
the pseudo-posterior sandwich approach is an important direction for
future work.

\paragraph{Conclusion.}
This paper introduces a Bayesian hierarchical hurdle beta-binomial
framework for survey-weighted bounded counts with structural zeros.
Applied to the 2019 NSECE, the model reveals a poverty reversal---community
poverty reduces access to infant and toddler care but increases
enrollment intensity among participating centers---that is
near-universal across states under hierarchical shrinkage (48 of 51)
and whose extensive margin accounts for approximately two-thirds of
the total effect.  The sandwich correction and DER diagnostic provide
a practical toolkit for population-level inference in Bayesian
hierarchical models fitted to complex survey data.  The companion R
package \textsf{hurdlebb}~\citep{hurdlebb2026} and replication
repository make the methodology immediately accessible to applied
researchers.

%% End-matter statements are in main.tex (after \input{section6}),
%% placed between the Discussion and References.
),
%% placed between the Discussion and References.
),
%% placed between the Discussion and References.
),
%% placed between the Discussion and References.
